\documentclass[11pt]{letter}

\usepackage[margin=0.55in]{geometry}
\usepackage{hyperref}

\begin{document}

\begin{letter}{BigGeo}

\opening{Dear Hiring Team,}

I am writing to apply for the Spatial AI Engineer position at BigGeo in Calgary. I am a Master of Science in Computer Science graduate from the University of Calgary, and my strongest experience is at the intersection this role describes: geospatial data infrastructure, machine learning, large spatial datasets, and practical engineering systems. Through my BigGeo/Mitacs research collaboration, I built C++ and Python software for scalable Earth observation data processing, storage, and retrieval, and I would be excited to continue contributing to BigGeo's Spatial Cloud work in a full-time engineering role.

My graduate research focused on making large satellite and terrain datasets easier to process, index, retrieve, and analyze. I designed a DGGS-based framework for real-time multiresolution management of spatiotemporal Earth observation data, developed tensor-based storage for direct access to arbitrary regions at multiple resolutions, and implemented wavelet and B-spline methods for spatial and temporal retrieval. This work included performance optimization in C++, automated preprocessing pipelines in Python, GDAL, Rasterio, SNAP, and satellite data APIs, and experiments at Canada-wide scale involving more than 1,100 GeoTIFF tiles and over 16,000 DGGS tensors.

I also bring hands-on machine learning experience with spatial data. In my DEM super-resolution work, I trained and evaluated PyTorch ResNet models to reconstruct 30m elevation data from 120m inputs, built a Streamlit and Rasterio pipeline to collect and preprocess mountainous terrain data, and evaluated model behavior using RMSE, slope, TRI, TPI, and wavelet-based metrics. That project strengthened the same habits BigGeo emphasizes: careful dataset construction, reliable evaluation, spatial reasoning, and practical trade-offs between accuracy, scalability, and usability.

Beyond research, I have backend engineering experience building Django REST APIs backed by PostgreSQL, implementing authentication and payment-related logic, and contributing to AI-enabled service workflows. As a teaching assistant for big data, data-at-scale, database, and deep learning courses, I also helped students debug pipelines, reason about cloud-based workflows, and implement PyTorch models. These experiences have made me comfortable moving between data engineering, machine learning, backend systems, and clear technical communication.

BigGeo is especially compelling to me because the role connects directly to the problems I have already been working on in Calgary: turning spatial and temporal data into systems that are fast, usable, and valuable for real decisions. I would welcome the opportunity to bring my geospatial research background, Python/C++ engineering experience, PyTorch model-development skills, and familiarity with BigGeo's domain to the Spatial AI Engineer role.

Thank you for considering my application. I look forward to the opportunity to discuss how my experience can support BigGeo's spatial AI and data infrastructure work.

\closing{Sincerely,

Amir Mirzai Golpayegani

\href{mailto:amirgolpaa24@gmail.com}{amirgolpaa24@gmail.com}}

\end{letter}

\end{document}
